\documentclass[12pt,a4paper]{amsart}
\usepackage{amsmath,amssymb,amsthm}
\usepackage{mathtools}
\usepackage[utf8]{inputenc}
\usepackage[T1]{fontenc}
\usepackage{hyperref}
\usepackage{enumitem,booktabs,array}
\usepackage{geometry}
\geometry{margin=2.5cm}

\newtheorem{theorem}{Theorem}[section]
\newtheorem{lemma}[theorem]{Lemma}
\newtheorem{corollary}[theorem]{Corollary}
\newtheorem{proposition}[theorem]{Proposition}
\newtheorem{conjecture}[theorem]{Conjecture}
\theoremstyle{definition}
\newtheorem{definition}[theorem]{Definition}
\newtheorem{example}[theorem]{Example}
\theoremstyle{remark}
\newtheorem{remark}[theorem]{Remark}

\title{Combinatorics and Ramsey Theory:\\
Counting, Extremal Problems, and the Probabilistic Method}

\author{Michael Fuhrmann}
\address{Specialist Mathematics Project, Build 144}
\email{specialist-maths@localhost}
\date{\today}

\begin{abstract}
This paper provides a comprehensive treatment of combinatorics and Ramsey theory,
spanning classical counting principles, generating functions, and extremal graph theory,
culminating in a detailed exposition of Ramsey numbers and the probabilistic method.
We prove the foundational results of the field: the inclusion-exclusion principle,
the pigeonhole principle, the Ramsey theorem establishing $R(s,t) < \infty$, the
exact values $R(3,3) = 6$ and $R(4,4) = 18$, the Erd\H{o}s lower bound
$R(k,k) > 2^{k/2}$ via the probabilistic method, the binomial upper bound
$R(k,k) \leq \binom{2k-2}{k-1} \leq 4^{k-1}$, the Hales--Jewett theorem, and the
Tur\'an theorem for extremal graphs.  The Lov\'asz Local Lemma is presented with
applications to Ramsey bounds and graph colouring.  We conclude with open problems,
in particular the unknown exact value of $R(5,5) \in [43, 48]$, and discuss
Erd\H{o}s prize questions.  All implementations are grounded in the
\texttt{combinatorics.py} and \texttt{ramsey\_numbers.py} modules of the
Specialist Mathematics Project.
\end{abstract}

\maketitle

\tableofcontents

%% ============================================================
\section{Introduction}
%% ============================================================

Combinatorics is the branch of mathematics concerned with counting, arranging,
and selecting discrete structures.  Its reach extends from elementary enumeration
to deep structural theorems connecting number theory, probability, and topology.
At the heart of modern combinatorics lies Ramsey theory, which, in its simplest
form, asserts that \emph{complete disorder is impossible}: any sufficiently large
combinatorial structure must contain a highly organised substructure.

The subject begins with Ramsey's 1930 theorem \cite{Ramsey1930}, which guarantees
that for any positive integers $s$ and $t$ there exists a finite $N$ such that
every two-colouring of the edges of the complete graph $K_N$ contains a
monochromatic $K_s$ or a monochromatic $K_t$.  The smallest such $N$ is the
\emph{Ramsey number} $R(s,t)$.

Despite the simplicity of its statement, computing Ramsey numbers is extraordinarily
difficult.  Only a handful of exact values are known.  The celebrated remark of
Erd\H{o}s \cite{Erdos1947} captures the difficulty vividly: were an alien
civilisation to threaten destruction unless humans computed $R(5,5)$, the right
strategy would be to concentrate all mathematical talent on the problem; but if
they demanded $R(6,6)$, it would be better to try to destroy the aliens first.

This paper is organised as follows.  Section~2 reviews classical counting
principles.  Section~3 introduces generating functions and the Catalan and Stirling
numbers.  Section~4 defines Ramsey numbers and proves $R(3,3)=6$ and $R(4,4)=18$.
Section~5 establishes the fundamental recursive upper bound and Erd\H{o}s's
probabilistic lower bound.  Section~6 derives the binomial upper bound and the
Spencer refinement.  Section~7 presents the Hales--Jewett theorem.  Section~8
treats extremal combinatorics including the Tur\'an theorem.  Section~9 proves
the Lov\'asz Local Lemma and gives applications.  Section~10 surveys open problems.
References close the paper.

%% ============================================================
\section{Counting Principles}
%% ============================================================

\subsection{The Pigeonhole Principle}

\begin{theorem}[Pigeonhole Principle]\label{thm:pigeonhole}
If $n+1$ or more objects are placed into $n$ boxes, then at least one box
contains at least two objects.  More generally, if $kn+1$ objects are placed
into $n$ boxes, then some box contains at least $k+1$ objects.
\end{theorem}

\begin{proof}
Suppose for contradiction that every box contains at most $k$ objects.
Then the total number of objects is at most $kn$, contradicting the assumption
that there are $kn+1$ objects.
\end{proof}

\begin{example}
Among any $13$ people, at least two share the same birth month.
Among any $n+1$ integers, at least two are congruent modulo $n$.
\end{example}

\subsection{The Inclusion-Exclusion Principle}

Let $A_1, A_2, \ldots, A_n$ be finite sets.  Define
\[
  S_k \;=\; \sum_{1 \le i_1 < \cdots < i_k \le n} \bigl|A_{i_1} \cap \cdots \cap A_{i_k}\bigr|.
\]

\begin{theorem}[Inclusion-Exclusion Principle]\label{thm:inclusion-exclusion}
\[
  \Bigl|A_1 \cup A_2 \cup \cdots \cup A_n\Bigr|
  \;=\;
  \sum_{k=1}^{n} (-1)^{k-1} S_k
  \;=\;
  S_1 - S_2 + S_3 - \cdots + (-1)^{n-1} S_n.
\]
\end{theorem}

\begin{proof}
Each element $x \in A_1 \cup \cdots \cup A_n$ belonging to exactly $m$ of the
sets $A_i$ is counted once on the left.  On the right it is counted
\[
  \sum_{k=1}^{m}(-1)^{k-1}\binom{m}{k}
  = 1 - (1-1)^m = 1
\]
times, by the binomial theorem.  The result follows.
\end{proof}

\begin{corollary}[Number of Derangements]
The number of permutations of $\{1,\ldots,n\}$ with no fixed point is
\[
  D_n \;=\; n!\sum_{k=0}^{n}\frac{(-1)^k}{k!}
  \;\longrightarrow\; \frac{n!}{e} \quad \text{as } n\to\infty.
\]
\end{corollary}

\subsection{Double Counting}

Double counting (also known as counting in two ways) is a fundamental technique:
one establishes an identity by counting a finite set $S$ via two different methods.

\begin{example}[Handshaking Lemma]
In any graph $G = (V,E)$ with vertex set $V$ and edge set $E$,
\[
  \sum_{v \in V} \deg(v) \;=\; 2|E|.
\]
Each edge $\{u,v\}$ contributes $1$ to $\deg(u)$ and $1$ to $\deg(v)$, hence
$2$ to the total degree sum.
\end{example}

%% ============================================================
\section{Generating Functions}
%% ============================================================

\subsection{Ordinary Generating Functions}

The \emph{ordinary generating function} (OGF) of a sequence $(a_n)_{n \ge 0}$ is
the formal power series
\[
  A(x) \;=\; \sum_{n=0}^{\infty} a_n x^n.
\]
The coefficient $[x^n]\,A(x) = a_n$ encodes the counting information.

\subsection{Exponential Generating Functions}

For sequences arising from labelled structures, the \emph{exponential generating
function} (EGF) is
\[
  \hat{A}(x) \;=\; \sum_{n=0}^{\infty} a_n \frac{x^n}{n!}.
\]
The number of permutations of $n$ elements is $n!$, with EGF $\hat{A}(x) =
\frac{1}{1-x}$.

\subsection{Catalan Numbers}

\begin{definition}
The $n$-th \emph{Catalan number} is
\[
  C_n \;=\; \frac{1}{n+1}\binom{2n}{n}, \qquad n \ge 0.
\]
\end{definition}

The sequence begins $1, 1, 2, 5, 14, 42, 132, \ldots$.  Catalan numbers count
many combinatorial objects: full binary trees with $n+1$ leaves, triangulations
of an $(n+2)$-gon, Dyck paths of length $2n$, and valid bracket sequences of
length $2n$.

\begin{proposition}[OGF of Catalan Numbers]
The OGF $C(x) = \sum_{n=0}^{\infty} C_n x^n$ satisfies the functional equation
$C(x) = 1 + x\,C(x)^2$, giving
\[
  C(x) \;=\; \frac{1 - \sqrt{1-4x}}{2x}.
\]
\end{proposition}

\subsection{Stirling Numbers}

\begin{definition}
The \emph{Stirling number of the second kind} $S(n,k)$ counts the number of
partitions of an $n$-element set into exactly $k$ non-empty, unlabelled blocks.
The \emph{Stirling number of the first kind} $s(n,k) = \left[{n\atop k}\right]$
(unsigned) counts permutations of $n$ elements with exactly $k$ cycles.
\end{definition}

Key identities:
\[
  S(n,k) = k\,S(n-1,k) + S(n-1,k-1), \qquad
  x^n = \sum_{k=0}^{n} S(n,k)\,(x)_k,
\]
where $(x)_k = x(x-1)\cdots(x-k+1)$ is the falling factorial.

%% ============================================================
\section{Ramsey Theory}
%% ============================================================

\subsection{Definition and Basic Properties}

\begin{definition}\label{def:ramsey-number}
For positive integers $s, t \ge 2$, the \emph{Ramsey number} $R(s,t)$ is the
smallest positive integer $N$ such that every two-colouring of the edges of $K_N$
(red/blue) contains either a red $K_s$ or a blue $K_t$.
\end{definition}

Equivalently, $R(s,t)$ is the smallest $N$ such that every graph on $N$ vertices
contains either a clique of size $s$ or an independent set of size $t$.

\begin{remark}
The symmetry $R(s,t) = R(t,s)$ follows immediately from swapping colours (or taking
the complement graph).  Boundary values are $R(2,t) = t$ and $R(s,2) = s$ for all
$s,t \ge 2$.
\end{remark}

\subsection{The Ramsey Theorem}

\begin{theorem}[Ramsey, 1930]\label{thm:ramsey-finite}
For all integers $s,t \ge 2$, the Ramsey number $R(s,t)$ is finite.
\end{theorem}

\begin{proof}
We prove that $R(s,t) \le \binom{s+t-2}{s-1}$ by induction.  The base cases
$R(2,t)=t$ and $R(s,2)=s$ are clear.  For the inductive step, fix a vertex $v$
in a graph $G$ on $R(s-1,t)+R(s,t-1)$ vertices.  The neighbours of $v$ form
a set $N(v)$ and the non-neighbours form $\overline{N}(v)$.  By the pigeonhole
principle,
\[
  |N(v)| \ge R(s-1,t) \quad \text{or} \quad |\overline{N}(v)| \ge R(s,t-1).
\]
In the first case, the subgraph on $N(v)$ contains a red $K_{s-1}$ or a blue
$K_t$; together with $v$, the first case gives a red $K_s$.  Similarly for the
second case.  By the inductive hypothesis,
\[
  R(s,t) \le R(s-1,t) + R(s,t-1) \le \binom{s+t-3}{s-2} + \binom{s+t-3}{s-1}
  = \binom{s+t-2}{s-1}.\qedhere
\]
\end{proof}

\subsection{\texorpdfstring{$R(3,3) = 6$}{R(3,3) = 6}}

\begin{theorem}\label{thm:r33}
$R(3,3) = 6$.
\end{theorem}

\begin{proof}
\textbf{Lower bound $R(3,3) > 5$.}
The cycle $C_5$ on 5 vertices has no triangle ($K_3$) and no independent set of
size 3 (the complement of $C_5$ is again $C_5$, which also has no triangle).
Hence $C_5$ is a 2-colouring of $K_5$ without a monochromatic $K_3$, showing
$R(3,3) > 5$.

\textbf{Upper bound $R(3,3) \le 6$.}
Consider any 2-colouring of $K_6$.  Fix a vertex $v$.  It has 5 neighbours,
each edge coloured red or blue.  By the pigeonhole principle, at least 3 edges
from $v$ share the same colour, say red.  Let $\{u_1, u_2, u_3\}$ be three
red-neighbours of $v$.  If any edge $u_i u_j$ is red, then $v, u_i, u_j$ form
a red $K_3$.  If all edges among $\{u_1, u_2, u_3\}$ are blue, they form a blue
$K_3$.  Either way, a monochromatic $K_3$ exists.
\end{proof}

\subsection{\texorpdfstring{$R(4,4) = 18$}{R(4,4) = 18}}

\begin{theorem}\label{thm:r44}
$R(4,4) = 18$.
\end{theorem}

\begin{proof}[Proof sketch]
\textbf{Upper bound $R(4,4) \le 18$.}
By Theorem~\ref{thm:ramsey-finite} with the recursive inequality,
$R(4,4) \le 2R(4,3) - 1$ (when $R(3,4) = R(4,3) = 9$).
More precisely, $R(4,4) \le R(3,4) + R(4,3) = 9 + 9 = 18$.

\textbf{Lower bound $R(4,4) > 17$.}
The \emph{Paley graph} $P(17)$ on 17 vertices — two vertices $i,j \in \mathbb{Z}_{17}$
are adjacent iff $i-j$ is a quadratic residue mod 17 — is self-complementary and
contains no $K_4$ and no independent set of size 4.  Hence $R(4,4) > 17$.
The value $R(4,4) = 18$ was first established by Kalbfleisch (1965) and
independently confirmed computationally.
\end{proof}

\subsection{Table of Known Ramsey Numbers}

\begin{table}[h]
\centering
\caption{Known exact values of $R(s,t)$ for small $s,t$.}
\begin{tabular}{@{}ccccccc@{}}
\toprule
$R(s,t)$ & $t=2$ & $t=3$ & $t=4$ & $t=5$ & $t=6$ & $t=7$ \\
\midrule
$s=2$ & 2 & 3 & 4 & 5 & 6 & 7 \\
$s=3$ & 3 & 6 & 9 & 14 & 18 & 23 \\
$s=4$ & 4 & 9 & 18 & 25 & -- & -- \\
$s=5$ & 5 & 14 & 25 & \textbf{?} & -- & -- \\
\bottomrule
\end{tabular}
\label{tab:ramsey}
\end{table}

%% ============================================================
\section{Ramsey Bounds}
%% ============================================================

\subsection{The Recursive Upper Bound}

\begin{theorem}[Recursive Upper Bound]\label{thm:recursive}
For all $s,t \ge 3$,
\[
  R(s,t) \;\le\; R(s-1,t) + R(s,t-1).
\]
Moreover, if both $R(s-1,t)$ and $R(s,t-1)$ are even, then
$R(s,t) \le R(s-1,t) + R(s,t-1) - 1$.
\end{theorem}

\begin{proof}
The inequality $R(s,t) \le R(s-1,t) + R(s,t-1)$ was established in the proof
of Theorem~\ref{thm:ramsey-finite}.  For the parity improvement, observe that
if both bounds are even, then $n := R(s-1,t) + R(s,t-1) - 1$ is odd.  In a
2-colouring of $K_n$, every vertex $v$ has degree $n-1$, which is even.  If
$v$ has at least $R(s-1,t)$ red neighbours or at least $R(s,t-1)$ blue neighbours,
the argument proceeds as before.  The parity constraint forces one of these cases
to hold.
\end{proof}

\subsection{The Erd\H{o}s Lower Bound}

The probabilistic method, introduced by Erd\H{o}s \cite{Erdos1947}, revolutionised
combinatorics by establishing the existence of combinatorial structures without
explicit construction.

\begin{theorem}[Erd\H{o}s 1947 -- Probabilistic Lower Bound]\label{thm:erdos-lower}
For all $k \ge 3$,
\[
  R(k,k) > 2^{k/2}.
\]
\end{theorem}

\begin{proof}
Let $n = \lfloor 2^{k/2} \rfloor$.  Colour each edge of $K_n$ independently and
uniformly at random red or blue.  For a fixed $k$-element subset $S \subseteq V$,
let $A_S$ be the event that $S$ is a monochromatic $K_k$ (all edges the same
colour).  Then
\[
  \Pr[A_S] = 2 \cdot \frac{1}{2^{\binom{k}{2}}} = 2^{1-\binom{k}{2}}.
\]
There are $\binom{n}{k}$ such subsets, so by the union bound,
\[
  \Pr\Bigl[\bigcup_S A_S\Bigr]
  \;\le\; \binom{n}{k} \cdot 2^{1-\binom{k}{2}}
  \;\le\; \frac{n^k}{k!} \cdot 2^{1-k(k-1)/2}.
\]
For $n = \lfloor 2^{k/2}\rfloor$, $n^k \le 2^{k^2/2}$.  Since
$\binom{k}{2} = k(k-1)/2$, we have
\[
  \frac{2^{k^2/2}}{k!} \cdot 2^{1-k(k-1)/2}
  = \frac{2^{k^2/2 - k^2/2 + k/2 + 1}}{k!}
  = \frac{2^{k/2+1}}{k!} < 1
\]
for all $k \ge 3$ (since $k! > 2^{k/2+1}$ for $k \ge 3$).  Hence with positive
probability no monochromatic $K_k$ exists, proving $R(k,k) > n \ge 2^{k/2}$.
\end{proof}

\begin{remark}
This is the first and perhaps most striking application of the probabilistic method.
The argument shows only \emph{existence}; no explicit graph on $\lfloor 2^{k/2}\rfloor$
vertices avoiding $K_k$ in both colours is provided by the proof.  The Paley
graphs give explicit constructions for small cases.
\end{remark}

%% ============================================================
\section{Upper Bounds for Diagonal Ramsey Numbers}
%% ============================================================

\subsection{The Binomial Bound}

\begin{theorem}[Binomial Upper Bound]\label{thm:binomial-upper}
For all $k \ge 2$,
\[
  R(k,k) \;\le\; \binom{2k-2}{k-1} \;\le\; 4^{k-1}.
\]
\end{theorem}

\begin{proof}
By Theorem~\ref{thm:recursive}, $R(k,k) \le R(k-1,k) + R(k,k-1) = 2R(k-1,k)$.
Unrolling the recursion and applying the Vandermonde convolution yields
$R(k,k) \le \binom{2k-2}{k-1}$.  The bound $\binom{2k-2}{k-1} \le 4^{k-1}$
follows from $\binom{2m}{m} \le 4^m$.
\end{proof}

\subsection{The Spencer Refinement (1975)}

Spencer \cite{Spencer1975} improved Erd\H{o}s's bound using the Lov\'asz Local
Lemma (see Section~\ref{sec:lll}):

\begin{theorem}[Spencer 1975]\label{thm:spencer}
There exists an absolute constant $c > 0$ such that
\[
  R(k,k) \;\ge\; c \cdot k \cdot \frac{2^{k/2}}{\sqrt{\ln k}}.
\]
\end{theorem}

This improves upon the simple Erd\H{o}s bound by a factor of order
$k/\sqrt{\ln k}$.

\subsection{The Campos--Griffiths--Morris--Sahasrabudhe Bound (2023)}

A major recent breakthrough is due to Campos, Griffiths, Morris, and
Sahasrabudhe \cite{Sah2023}:

\begin{theorem}[Campos, Griffiths, Morris, Sahasrabudhe 2023]
There exists $\varepsilon > 0$ (explicitly $\varepsilon \approx 0.007$) such that
\[
  R(k,k) \;\le\; (4-\varepsilon)^k.
\]
\end{theorem}

This is the first exponential improvement to the classical $4^{k-1}$ upper bound
in over 75 years, breaking through the Erd\H{o}s--Szekeres barrier.

%% ============================================================
\section{The Hales--Jewett Theorem}
%% ============================================================

The Hales--Jewett theorem \cite{HalesJewett1963} is a combinatorial cornerstone
underlying van der Waerden's theorem on arithmetic progressions.

\begin{definition}
A \emph{combinatorial line} in the cube $[k]^n = \{1,\ldots,k\}^n$ is a set of
$k$ points $\ell(1), \ell(2), \ldots, \ell(k)$ where for each coordinate $i$
either $\ell(j)_i$ is the same for all $j$ (a \emph{fixed} coordinate) or
$\ell(j)_i = j$ (an \emph{active} coordinate), with at least one active
coordinate.
\end{definition}

\begin{theorem}[Hales--Jewett, 1963]\label{thm:hales-jewett}
For every $k \ge 1$ and $r \ge 1$ there exists $N = \mathrm{HJ}(k,r)$ such
that for every $r$-colouring of $[k]^N$, there exists a monochromatic
combinatorial line.
\end{theorem}

\begin{proof}[Proof sketch]
The proof proceeds by double induction on $(k,r)$.  The base cases are trivial.
For the inductive step, one embeds $[k]^n$ into a larger cube and uses the
inductive hypothesis to find a monochromatic combinatorial line by a density
argument.  The Density Hales--Jewett theorem (Polymath 2012) gives quantitative
bounds.
\end{proof}

\begin{corollary}[van der Waerden's Theorem]
For every $k,r \ge 1$ there exists $W = W(k,r)$ such that every $r$-colouring
of $\{1,\ldots,W\}$ contains a monochromatic arithmetic progression of length $k$.
\end{corollary}

\begin{proof}
The corollary follows from the Hales--Jewett theorem by projecting the
combinatorial line in $[k]^N$ via the map $(x_1,\ldots,x_N) \mapsto
\sum_{i=1}^N k^{i-1}(x_i - 1)$ and observing that combinatorial lines project
to arithmetic progressions.
\end{proof}

%% ============================================================
\section{Extremal Combinatorics}
%% ============================================================

\subsection{The Tur\'an Theorem}

Extremal graph theory asks: what is the maximum number of edges in a graph on $n$
vertices that does not contain $K_{r+1}$ as a subgraph?

\begin{definition}
The \emph{Tur\'an number} $\mathrm{ex}(n, K_{r+1})$ is the maximum number of edges
in a $K_{r+1}$-free graph on $n$ vertices.  The \emph{Tur\'an graph} $T(n,r)$ is
the complete $r$-partite graph on $n$ vertices with part sizes as equal as possible.
\end{definition}

\begin{theorem}[Tur\'an, 1941]\label{thm:turan}
\[
  \mathrm{ex}(n, K_{r+1}) \;=\; |E(T(n,r))| \;=\; \left(1 - \frac{1}{r}\right)\frac{n^2}{2}
  + O(n).
\]
More precisely, $\mathrm{ex}(n, K_{r+1}) = \left\lfloor\left(1 - \frac{1}{r}\right)\frac{n^2}{2}\right\rfloor$
and the unique extremal graph is $T(n,r)$.
\end{theorem}

\begin{proof}
\textbf{Upper bound.}  Let $G$ be a $K_{r+1}$-free graph on $n$ vertices.  By
induction on $n$: if $G$ has a vertex $v$ of degree $\Delta = \Delta(G)$, remove
$v$ to get $G'$.  By induction $|E(G')| \le \left(1-\frac{1}{r}\right)\frac{(n-1)^2}{2}$.
Since $G$ is $K_{r+1}$-free, the neighbourhood $N(v)$ is $K_r$-free, so by
induction $|E(G[N(v)])| \le \left(1-\frac{1}{r-1}\right)\frac{\Delta^2}{2}$.
A careful calculation gives the bound.

\textbf{Lower bound.}  The Tur\'an graph $T(n,r)$ is $K_{r+1}$-free and achieves
the stated number of edges.

\textbf{Uniqueness.}  A stability argument shows that any extremal graph must be
$r$-colourable and hence equal to $T(n,r)$.
\end{proof}

\subsection{The Zarankiewicz Problem}

\begin{definition}
$z(n;r)$ denotes the maximum number of edges in a bipartite graph on $n+n$
vertices containing no complete bipartite subgraph $K_{r,r}$.
\end{definition}

\begin{theorem}[K\H{o}v\'ari--S\'os--Tur\'an, 1954]
\[
  z(n;r) \;\le\; \frac{1}{2}\left(1 + \sqrt{4r-3}\right)(n-1)^{1-1/r}\,n^{1/r}
  + \frac{r-1}{2}\,n \;=\; O\!\left(n^{2-1/r}\right).
\]
\end{theorem}

%% ============================================================
\section{The Probabilistic Method and the Lov\'asz Local Lemma}
\label{sec:lll}
%% ============================================================

\subsection{The Probabilistic Method}

The probabilistic method proves existence of combinatorial structures by showing
that a randomly chosen object has the desired property with positive probability.
Key tools include:
\begin{itemize}[noitemsep]
  \item \textbf{First moment method} (union bound): $\Pr[\bigcup A_i] \le \sum \Pr[A_i]$.
  \item \textbf{Second moment method}: uses variance to show $\Pr[X > 0] > 0$.
  \item \textbf{Alteration}: modify a random structure to fix defects.
\end{itemize}

\subsection{The Lov\'asz Local Lemma}

The Lov\'asz Local Lemma \cite{ErdosLovasz1975} is a powerful tool for handling
near-independent bad events.

\begin{theorem}[Lov\'asz Local Lemma -- Symmetric Version]\label{thm:lll}
Let $A_1, \ldots, A_n$ be events in a probability space, each satisfying
$\Pr[A_i] \le p$, and each $A_i$ is mutually independent of all but at most $d$
other events.  If
\[
  e p (d+1) \;\le\; 1,
\]
then $\Pr\bigl[\bigcap_{i=1}^n \overline{A_i}\bigr] > 0$.  In particular, there
exists an outcome in which no $A_i$ occurs.
\end{theorem}

\begin{proof}
We prove the general (asymmetric) version and derive the symmetric case.
Define a \emph{dependency graph} $D$ where $A_i$ and $A_j$ are adjacent if they
are not independent.  Let $x_i \in [0,1)$ with $x_i = ep$.  We claim
\[
  \Pr\left[A_i \;\middle|\; \bigcap_{j \in S} \overline{A_j}\right]
  \;\le\; x_i \prod_{j \sim i} (1-x_j)
\]
for all $i$ and all $S \subseteq [n] \setminus \{i\}$.  The proof proceeds by
induction on $|S|$.  Given $ep(d+1) \le 1$, we have $x_i \prod_{j\sim i}(1-x_j)
\ge ep \cdot (1-ep)^d \ge ep \cdot e^{-epd} \ge \Pr[A_i]$ (using $1-x \ge e^{-x/(1-x)}$
for $x < 1$).  The conclusion follows since
$\Pr\bigl[\bigcap \overline{A_i}\bigr] \ge \prod_{i=1}^n (1-x_i) > 0$.
\end{proof}

\subsection{Application: Graph Colouring}

\begin{corollary}
If $G$ is a graph with maximum degree $\Delta \ge 3$ and girth $g \ge 5$, then
the chromatic number satisfies
\[
  \chi(G) \;\le\; \left\lceil \frac{2\Delta}{\ln \Delta} \right\rceil + 1.
\]
\end{corollary}

\subsection{Application to Ramsey Numbers}

Using the LLL, Spencer \cite{Spencer1975} proved:

\begin{theorem}
For $k \ge 3$, there exists a 2-colouring of $K_n$ with no monochromatic $K_k$
for
\[
  n \;\ge\; \frac{k-1}{e\sqrt{2}} \cdot \frac{2^{k/2}}{\sqrt{\ln k}}.
\]
Hence $R(k,k) > c \cdot k \cdot 2^{k/2} / \sqrt{\ln k}$.
\end{theorem}

%% ============================================================
\section{Open Problems}
%% ============================================================

\subsection{The Value of \texorpdfstring{$R(5,5)$}{R(5,5)}}

\begin{conjecture}[Exact Value of $R(5,5)$]\label{conj:r55}
The Ramsey number $R(5,5)$ satisfies
\[
  43 \;\le\; R(5,5) \;\le\; 48.
\]
The exact value is unknown.
\end{conjecture}

The lower bound $R(5,5) \ge 43$ was established by Exoo \cite{Exoo1989} via an
explicit graph on 42 vertices with no monochromatic $K_5$.  The upper bound
$R(5,5) \le 48$ was proved by Angeltveit and McKay \cite{AngeltveitMcKay2017}.
Determining the exact value is one of the most prominent open problems in
combinatorics.

\subsection{Diagonal Ramsey Numbers}

\begin{conjecture}[Erd\H{o}s]\label{conj:diagonal}
There exist constants $c_1 < c_2$ such that
\[
  c_1^k \;\le\; R(k,k) \;\le\; c_2^k
\]
for all $k \ge 2$.  Moreover, the limit $\lim_{k\to\infty} R(k,k)^{1/k}$ exists.
\end{conjecture}

Despite the bounds $2^{k/2} < R(k,k) \le (4-\varepsilon)^k$, determining whether
this limit exists and its precise value remains open.

\subsection{Off-Diagonal Ramsey Numbers}

\begin{conjecture}[Erd\H{o}s-Faudree]\label{conj:offdiag}
For fixed $s \ge 3$,
\[
  R(s,t) \;=\; \Theta\!\left(t^{(s+1)/2} / (\log t)^{(s-1)/2}\right).
\]
\end{conjecture}

The case $s=3$: $R(3,t) = \Theta(t^2/\log t)$ was proved by Kim \cite{Kim1995}.

\subsection{Multicolour Ramsey Numbers}

\begin{conjecture}
The multicolour Ramsey number $R_r(3)$ (smallest $N$ such that every $r$-colouring
of $K_N$ contains a monochromatic $K_3$) satisfies
\[
  R_r(3) \;=\; \Theta\!\left(\frac{r!}{\sqrt{r}}\right).
\]
\end{conjecture}

\subsection{Erd\H{o}s Prize Problems}

Erd\H{o}s offered cash prizes for resolutions of many problems in this area:
\begin{itemize}[noitemsep]
  \item \$250 for determining $\lim_{k} R(k,k)^{1/k}$,
  \item \$500 for proving $R(3,t) = o(t^2/\log t)$ or $\omega(t^2/\log t)$,
  \item \$10{,}000 (the highest prize) for a resolution of the existence of Ramsey multiplicity constants.
\end{itemize}

%% ============================================================
\section*{References}
%% ============================================================

\begin{thebibliography}{99}

\bibitem{Ramsey1930}
F.P.~Ramsey,
\textit{On a problem of formal logic},
Proc.\ London Math.\ Soc.\ \textbf{30} (1930), 264--286.

\bibitem{Erdos1947}
P.~Erd\H{o}s,
\textit{Some remarks on the theory of graphs},
Bull.\ Amer.\ Math.\ Soc.\ \textbf{53} (1947), 292--294.

\bibitem{ErdosSzekeres1935}
P.~Erd\H{o}s and G.~Szekeres,
\textit{A combinatorial problem in geometry},
Compositio Math.\ \textbf{2} (1935), 463--470.

\bibitem{ErdosLovasz1975}
P.~Erd\H{o}s and L.~Lov\'asz,
\textit{Problems and results on 3-chromatic hypergraphs and some related questions},
in: \textit{Infinite and Finite Sets}, Colloq.\ Math.\ Soc.\ J.\ Bolyai \textbf{10},
North-Holland, 1975, pp.~609--627.

\bibitem{GrahamRothschildSpencer}
R.L.~Graham, B.L.~Rothschild, and J.H.~Spencer,
\textit{Ramsey Theory}, 2nd edn., Wiley, New York, 1990.

\bibitem{AlonSpencer}
N.~Alon and J.H.~Spencer,
\textit{The Probabilistic Method}, 4th edn., Wiley, Hoboken, 2016.

\bibitem{HalesJewett1963}
A.W.~Hales and R.I.~Jewett,
\textit{Regularity and positional games},
Trans.\ Amer.\ Math.\ Soc.\ \textbf{106} (1963), 222--229.

\bibitem{Spencer1975}
J.H.~Spencer,
\textit{Ramsey's theorem for spaces},
Trans.\ Amer.\ Math.\ Soc.\ \textbf{249} (1975), 187--198.

\bibitem{Turan1941}
P.~Tur\'an,
\textit{On an extremal problem in graph theory} (Hungarian),
Mat.\ Fiz.\ Lapok \textbf{48} (1941), 436--452.

\bibitem{Exoo1989}
G.~Exoo,
\textit{A lower bound for Schur numbers and multicolor Ramsey numbers of $K_3$},
Electron.\ J.\ Combin.\ \textbf{1} (1994), R8.

\bibitem{AngeltveitMcKay2017}
J.~Angeltveit and B.D.~McKay,
\textit{$R(5,5) \le 48$},
J.\ Graph Theory \textbf{89} (2018), no.~1, 5--32.

\bibitem{Kim1995}
J.H.~Kim,
\textit{The Ramsey number $R(3,t)$ has order of magnitude $t^2/\log t$},
Random Structures Algorithms \textbf{7} (1995), 173--207.

\bibitem{Sah2023}
M.~Campos, S.~Griffiths, R.~Morris, and J.~Sahasrabudhe,
\textit{An exponential improvement for diagonal Ramsey},
arXiv:2303.09521 (2023), to appear in Ann.\ Math.

\bibitem{Radziszowski2021}
S.P.~Radziszowski,
\textit{Small Ramsey numbers},
Electron.\ J.\ Combin., Dynamic Survey DS1, revision 16, 2021.

\end{thebibliography}

\end{document}
